%% =====================================================================
%%  T-19-02  Unequal Exchange
%%  -------------------------------------------------------------------
%%  One idea per file. Core is mandatory. Slots in canonical order.
%%  Max 700 words. No layout commands. No filler. New source material
%%  for this entry goes in sources/<BOOK>/T-19-02.tex -- never by
%%  rewriting this file (docs/RULES.md R0.1).
%% =====================================================================
%% @kind: topic
%% @id: T-19-02
%% @title: Unequal Exchange
%% @subtitle: The repayment need not resemble the debt
%% @chapter: c19
%% @order: 20
%% @status: stable
%% @sources: B4
%% @updated: 2026-09-19

\topic{T-19-02}{Unequal Exchange}{The repayment need not resemble the debt}

\begin{core}
A small benefit can be traded for a much larger one, and the trade usually goes through. The pressure generated is to repay, not to repay equally, and nothing in the mechanism sets an exchange rate.
\end{core}
\begin{mechanism}
The account carries no stated amount and no stated currency, so the debtor settles it by reference to a feeling rather than to an arithmetic of equivalence. That leaves the size of the repayment to be proposed by the creditor, who has every reason to propose large.

Two forces reinforce the asymmetry. The first is discomfort: an unsettled account is unpleasant to hold, and the quickest way to be rid of the feeling is to agree to whatever is asked, which makes the yes a relief rather than a decision. The second is exposure: not reciprocating is publicly marked, so the debtor is not merely uncomfortable but visibly so, and even a lopsided repayment clears the marking. Because both forces respond to the fact of the debt and not to its size, the disproportion is never felt as a reason to refuse.
\end{mechanism}
\begin{tells}
  \item The request is far larger than the benefit, and the disproportion goes unremarked.
  \item You experience owing a decision rather than owing a comparable favour.
  \item Agreeing relieves a feeling more than it achieves an outcome.
  \item The benefit was chosen by the other party and costs them little to give.
  \item No price was ever attached to the benefit, and attaching one now feels rude.
\end{tells}
\begin{moves}
  \item Keep the opening benefit small, cheap for you to give and visibly personal; the exchange rate is set by the feeling, not by the value.
  \item Ask for the larger thing soon afterwards, and in a different currency, so that no comparison is available to the other party.
  \item Leave the benefit unpriced. A benefit that has been costed has already been repaid.
\end{moves}
\begin{counters}
  \item Insist on the arithmetic. Establish what the benefit was worth, repay in kind and promptly, and treat the account as closed.
  \item Name the disproportion plainly and without accusation: this is considerably larger than what you did for me.
  \item Delay. The discomfort driving the unequal exchange decays faster than the obligation itself.
  \item Choose the repayment currency rather than accepting the one offered.
  \item Ask what the benefit cost the giver. A cheap gift demanding an expensive return is a transaction, not a kindness, and can be treated as one.
\end{counters}
\begin{cost}
Pushed too far the technique brands its user: a person known for trading small kindnesses into large concessions finds the marking that makes reciprocity work running against them, and later benefits are refused at the door. It also fails outright where the account is not meant to be settled at all --- family, long colleagues --- because proposing an exchange rate in such a relationship is itself the offence, and the offence outlasts the concession.
\end{cost}
\begin{field}
A free sample is not a gift of product. It is the opening of an account whose intended repayment is a purchase, and whose size is set at the till rather than at the counter.
\end{field}

\sources{B4}
\seesources
\seealso{T-19-01, T-04-02, T-25-01, T-24-01}
